\documentclass[a4paper, oneside, 11pt]{report}
    \usepackage[table]{xcolor}
    \usepackage[T1]{fontenc}
    \usepackage[utf8]{inputenc}
    \usepackage[nswissgerman, english]{babel}
    \usepackage[hyphens]{url}
    \usepackage{graphicx}
    \usepackage{subfig}
    \usepackage{vhistory}
    \usepackage{float}
    \usepackage{pdfpages}
    \usepackage{tcolorbox}
    \usepackage{xcolor}
    \usepackage{nameref}
    \usepackage{makecell} % for line breaks in tables
    \usepackage{longtable}
    \usepackage{titlesec}

    
    % Style des URL
    \urlstyle{sf}
    % URL-Umbrüche
    \appto\UrlBreaks{\do\a\do\b\do\c\do\d\do\e\do\f\do\g\do\h\do\i\do\j
\do\k\do\l\do\m\do\n\do\o\do\p\do\q\do\r\do\s\do\t\do\u\do\v\do\w
\do\x\do\y\do\z}

    % Seitenränder
    \usepackage{geometry}
    \geometry{
        a4paper,
        left=20mm,
        top=30mm,
        right=20mm,
        bottom=30mm
    }

    % Jede Überschrift 1 auf neuer Seite
    \let\stdsection\section
    \renewcommand\section{\clearpage\stdsection}

    % Header and footer
    \usepackage{fancyhdr}
    \pagestyle{fancy}
    \fancyhf{}
    \lhead{\small \TITLE \\ \vspace{0.5mm} \normalsize \nouppercase\rightmark \vspace{0.0cm}}
    \rhead{
        \begin{picture}
            (0,0) \put(-100,0){\includegraphics[width=0.2\linewidth]{./assets/logo/hsr.png}}
        \end{picture}}
    \cfoot{\thepage}

    % Multicolomns und Multirows
    \usepackage{multicol}
    \usepackage{multirow}
    \setlength{\multicolsep}{2.0pt plus 2.0pt minus 1.5pt}% 50% of original values (above/below multicols)

    % Chapter ohne Nummerierung, Eintrag in Inhaltsverzeichnis
    \newcommand{\kapitel}[1]{
        \stepcounter{chapter}\chapter*{#1}
        % \addcontentsline{toc}{chapter}{#1}
        \markboth{\arabic{chapter} #1}{\arabic{chapter} #1}
        
    }

    % Section ohne Nummerierung, Eintrag in Inhaltsverzeichnis
    \newcommand{\sectionroman}[1]{
        \stepcounter{section}\section*{#1}
        \addcontentsline{toc}{section}{#1}
        \markboth{\arabic{section} #1}{\arabic{section} #1}
    }

    % Section ohne Nummerierung, Eintrag in Inhaltsverzeichnis
    \newcommand{\subsectionroman}[1]{
        \stepcounter{subsection}\subsection*{#1}
        \addcontentsline{toc}{subsection}{#1}
        \markboth{\arabic{subsection} #1}{\arabic{subsection} #1}
    }

    % Anpassung der Inhaltsverzeichnis-Tiefe, beginnend bei section
    \renewcommand{\partname}{}
    \renewcommand{\thesection}{\arabic{section}}
    \setcounter{secnumdepth}{3}
    \setcounter{tocdepth}{3}

    % Code Listings
    \usepackage{listings}
    \usepackage{color}
    \usepackage{beramono}


    \definecolor{bluekeywords}{rgb}{0,0,1}
    \definecolor{greencomments}{rgb}{0,0.5,0}
    \definecolor{redstrings}{rgb}{0.64,0.08,0.08}
    \definecolor{xmlcomments}{rgb}{0.5,0.5,0.5}
    \definecolor{types}{rgb}{0.17,0.57,0.68}

    \lstdefinestyle{visual-studio-style}{
        language=[Sharp]C,
        columns=flexible,
        showstringspaces=false,
        basicstyle=\footnotesize\ttfamily, 
        commentstyle=\color{greencomments},
        morekeywords={partial, var, value, get, set},
        keywordstyle=\bfseries\color{bluekeywords},
        stringstyle=\color{redstrings},
        breaklines=true,
        breakatwhitespace=true,
        tabsize=4,
        numbers=left,
        numberstyle=\tiny\color{black},
        frame=lines,
        showspaces=false,
        showtabs=false,
        escapeinside={£}{£},
    }

    \definecolor{Maroon}{rgb}{0.5,0,0}
    \definecolor{darkgreen}{rgb}{0,0.5,0}
    
    \lstdefinelanguage{xml}
    {
        morekeywords={id, layout_width, layout_height, name, layout},
        alsoletter=-,
        morestring=[b]",
        stringstyle=\color[rgb]{0,0,1},
        morecomment=[s]{<?}{?>},
        morecomment=[s]{<!--}{-->},
        morecomment=[s]{<!}{>},
        commentstyle=\color{darkgreen},
        moredelim=[s][\color{black}]{![}{]]},
        moredelim=*[s][\color{Maroon}]{<}{>},
        keywordstyle=\color{red}
    }

    \definecolor{violetkeywords}{RGB}{190,119,227}
    \lstdefinelanguage{Powershell}
    {
        morecomment=[l]{\#},
        % morekeywords={foreach, if, elseif, in, continue, return, switch, try, catch},
        % keywordstyle=\bfseries\color{violetkeywords}
    }
    
    \lstset{style=visual-studio-style}
    \usepackage[hidelinks]{hyperref}
